%=============================================================================
\section{集群资源概览}
%=============================================================================

\subsection{分区（Partition）}

SLURM 把计算节点分成几组，叫做``分区''。不同分区有不同的用途和限制：

\begin{table}[h]
\centering
\begin{tabular}{llp{7cm}}
\toprule
分区 & 时间上限 & 说明 \\
\midrule
\texttt{batch} & 2 天 & 默认分区，仅 CPU。不允许请求 GPU。 \\
\texttt{gpu} & 2 天 & GPU 分区。必须用 \texttt{--gres} 请求 GPU，不允许只跑 CPU 任务。 \\
\texttt{preempt} & 2 天 & 资源最多（包含教授贡献的节点），但作业可能被节点所有者\textbf{抢占}——你的作业会被直接 kill。 \\
\bottomrule
\end{tabular}
\end{table}

\subsection{每用户资源限制}

\begin{table}[h]
\centering
\begin{tabular}{lccc}
\toprule
分区组合 & CPU 核数上限 & 内存上限 & GPU 上限 \\
\midrule
\texttt{batch} + \texttt{gpu} & 250 & 5,000 GB & 10 \\
\texttt{preempt} & 1,000 & 8,000 GB & 20 \\
\bottomrule
\end{tabular}
\end{table}

\begin{tipbox}
可以同时指定多个分区做 fallback：\texttt{\#SBATCH -p gpu,preempt}。这样如果 \texttt{gpu} 分区没资源，会自动尝试 \texttt{preempt}。
\end{tipbox}

\subsection{可用 GPU}

集群有多种型号的 GPU：

\begin{table}[h]
\centering
\begin{tabular}{lcl}
\toprule
GPU 型号 & 显存 & 适用场景 \\
\midrule
T4 & 16 GB & 小模型推理、轻量训练 \\
L40 / L40s & 48 GB & 中等规模训练 \\
A100-40G & 40 GB & 大规模训练 \\
A100-80G & 80 GB & 大模型训练 \\
H200 & 141 GB & 最大规模任务 \\
\bottomrule
\end{tabular}
\end{table}

请求特定型号的 GPU：

\begin{lstlisting}[language=slurm]
#SBATCH --gres=gpu:a100:1       # 请求 1 块 A100（不区分 40G/80G）
#SBATCH --constraint="a100-80G" # 加上这行可以指定 80G 版本
\end{lstlisting}

\begin{warnbox}
除非你的框架支持多 GPU 并行，否则\textbf{不要请求多块 GPU}——你只会浪费资源、延长排队时间。
\end{warnbox}

\subsection{怎样查看集群当前的资源占用情况？}

在提交作业之前，你可能想知道：现在哪个分区有空闲资源？GPU 还剩多少？这可以帮你决定选哪个分区、请求什么类型的 GPU，避免排长队。

\subsubsection{方法一：hpctools（推荐，最直观）}

Tufts 提供了一个叫 \texttt{hpctools} 的辅助工具，用交互式菜单帮你查看各种信息：

\begin{lstlisting}[language=bash]
module load hpctools
hpctools
\end{lstlisting}

运行后会显示一个菜单，输入编号选择功能：

\begin{table}[h]
\centering
\begin{tabular}{cl}
\toprule
选项 & 功能 \\
\midrule
1 & 查看各分区可用的 CPU 核数和内存 \\
2 & 查看各分区可用的 GPU 资源 \\
3 & 查看某日期之后你的已完成作业历史 \\
4 & 查看你当前运行的作业详情 \\
5 & 查看 Research 存储的配额和用量 \\
6 & 查看指定目录的详细空间占用 \\
7 & 查看 Home 目录的空间占用明细 \\
\bottomrule
\end{tabular}
\end{table}

比如你想在提交 GPU 作业前看看还有多少 GPU 可用，选 \texttt{2} 就能看到各分区的 GPU 空闲情况。想检查自己的存储快满了没有，选 \texttt{5} 或 \texttt{7}。

\subsubsection{方法二：sinfo（命令行快速查看）}

如果你不想进交互式菜单，\texttt{sinfo} 可以一行命令看到集群状态：

\begin{lstlisting}[language=bash]
# 查看所有分区的节点状态概览
sinfo

# 只看 gpu 分区，显示空闲/占用/总计
sinfo -p gpu -o "%P %a %D %C"
\end{lstlisting}

输出中 \texttt{\%C} 格式为 \texttt{已分配/空闲/其他/总计} 的 CPU 核数。

查看 GPU 分区中各节点的具体 GPU 空闲情况：

\begin{lstlisting}[language=bash]
# 显示每个节点的 GPU 分配情况
sinfo -p gpu -O NodeList,Gres,GresUsed,CPUsState,FreeMem
\end{lstlisting}

\subsubsection{方法三：OnDemand 网页界面}

如果你偏好图形化界面，OnDemand 提供了两个有用的页面：

\begin{itemize}[nosep]
  \item \textbf{Clusters $\rightarrow$ System Status}：查看整个集群各分区的节点状态
  \item \textbf{Jobs $\rightarrow$ Active Jobs}：查看你自己的作业运行情况，点击单个作业可以看详情，已完成的作业还能在 Resources and I/O 标签下看到 CPU 和内存的效率
\end{itemize}

\subsubsection{查看正在运行的作业的实时资源使用}

如果你的作业已经在跑了，想看它实际占用了多少 CPU 和 GPU：

\begin{lstlisting}[language=bash]
# SSH 到作业所在的节点（先用 squeue --me 找到节点名）
ssh pax006

# 查看你的进程的 CPU 和内存使用
top -u $USER

# 查看 GPU 使用情况（在有 GPU 的节点上）
nvidia-smi
\end{lstlisting}

\texttt{nvidia-smi} 会显示每块 GPU 的算力利用率、显存占用、以及正在使用 GPU 的进程。如果你发现 GPU 利用率很低（比如低于 30\%），可能说明你的程序瓶颈在数据加载而不是计算——可以尝试增加 \texttt{--cpus-per-task} 来加速数据预处理。
